\documentclass[11pt]{article}
\usepackage{explorer}
\newcommand{\ExplorerLanguage}{EN}
\newcommand{\ExplorerNumber}{07}
\newcommand{\ExplorerKind}{Worked solutions}
\begin{document}
\ExplorerTitle{High school · probability}{Reversing a conditional probability}
\ExplorerSection{Exercise data}
For exercises 1 and 2: choose urn A with probability \(\frac{1}{3}\) and urn B with probability \(\frac{2}{3}\). A contains 3 red balls and 1 blue; B contains 1 red and 3 blue. Then draw a ball uniformly from the chosen urn.\par
\ExplorerSection{Exercise 1}
In the example, find P(A ∩ blue).\par
\((\frac{1}{3})(\frac{1}{4}) = \frac{1}{12}\).\par
\ExplorerSection{Exercise 2}
Find \(P(A|\mathrm{blue})\).\par
\(P(\mathrm{blue}) = \frac{7}{12}\). Hence \(P(A|\mathrm{blue}) = \frac{\frac{1}{12}}{\frac{7}{12}} = \frac{1}{7}\).\par
\ExplorerSection{Exercise 3}
In a class, 12 of 30 students play chess; 9 of those enjoy puzzles. Find \(P(\mathrm{puzzles}|\mathrm{chess})\). Can \(P(\mathrm{chess}|\mathrm{puzzles})\) be deduced?\par
The first probability is \(\frac{9}{12} = \frac{3}{4}\). The total number who enjoy puzzles is needed for the second.\par
\ExplorerSection{Exercise 4}
Events A and R satisfy \(P(A)=0.4, P(R|A)=0.5\) and \(P(R|\mathrm{not} A)=0.5\). Find \(P(A|R)\).\par
\(P(R)=0.5\) and \(P(A \cap  R)=0.2\), so \(P(A|R)=0.4\). Here observing R does not change the probability of A.\par
\ExplorerSection{Problem: complete reasoning}
Choose A with probability \(\frac{1}{3}\) and B with probability \(\frac{2}{3}\). Urn A contains 3 red balls and 1 blue; B contains 1 red and 3 blue. After observing one red ball, the probability that A was chosen is \(\frac{3}{5}\).\par
Use the two urns from the example. Choose an urn once, then draw twice with replacement from that same urn. Find the probability of two red balls, then the probability that A was chosen given two red balls. Compare with observing one red ball and justify the products used.\par
Conditional on the selected urn, draws with replacement are independent. \(P(RR) = (\frac{1}{3})(\frac{3}{4})^{2} + (\frac{2}{3})(\frac{1}{4})^{2} = \frac{9}{48} + \frac{2}{48} = \frac{11}{48}\). Thus \(P(A|RR) = \frac{\frac{9}{48}}{\frac{11}{48}} = \frac{9}{11}\), greater than \(\frac{3}{5}\). The draws are not independent unconditionally: they share the same hidden urn.\par
\end{document}
